\setcounter{section}{0}
\setcounter{subsection}{0}
\setcounter{subsubsection}{0}

\clearpage
\begin{center}
  {\Large\bfseries Take-Home Lab}\\[0.5em]
\end{center}
\vspace{1em}

\section{Take-Home Portion: Empirical Evaluation of $N$-Version Programming}

In this take-home component, you will extend the in-class experiment in two directions:

\begin{itemize}
  \item Develop the general $k$-out-of-$N$ reliability model under the assumption of independence.
  \item Empirically evaluate whether LLM-generated implementations conform to that model.
\end{itemize}

\subsection{Theory: Derivation of the $k$-out-of-$N$ Model}

Suppose each version fails independently with probability $p$ on a given demand.
Let $q = 1 - p$ denote the probability of success.

Let $X$ be the number of successful versions out of $N$.
Under the independence assumption, $X$ follows a binomial distribution:

\[
X \sim \text{Binomial}(N, q).
\]

For a $k$-out-of-$N$ voting architecture, the system is correct if at least $k$ versions are correct.
Equivalently, the system fails if fewer than $k$ versions are correct.

\textbf{Your task:}
Derive an expression for the system failure probability $P(\text{system failure})$ for a general $k$-out-of-$N$ voter.
Your final answer should be written in summation form using binomial coefficients (a closed-form simplification is \emph{not} required).
Clearly indicate where and how the independence assumption is used in your derivation.

Your derivation should clearly specify:

\begin{itemize}
  \item The definition of the random variable $X$ (what does it count?),
  \item The probability that exactly $i$ versions are correct,
  \item The condition under which the $k$-out-of-$N$ voter fails,
  \item The resulting summation expression for $P(\text{system failure})$.
\end{itemize}

\subsection{Evaluation: Empirical Comparison to the Independence Model}

In the in-class portion, you initiated $N$-version development using LLM-generated implementations, created manually.
For the take-home component, you will scale this up.

\subsubsection{Step 1: Generate Versions}

Generate at least $N = 100$ distinct implementations of the provided specification.

You may use your choice of:

\begin{itemize}
  \item LLM-based agents,
  \item Interactive LLM sessions,
  \item Multiple fresh sessions with varied prompts,
  \item Any system available to you (\eg Gemini Pro, Microsoft Copilot, open-weight models).
\end{itemize}

If you have not used these types of systems before, here are some options with references to tutorials on how to get started:

\begin{enumerate}
    \item CoPilot Skills on VS Code\cite{copilot-documentation}. Pay close attention to the examples on the drop-down for this.

    \item A programmatic agent system such as python Transformers~\cite{transformers-tutorial}. This may allow for more control than the one above, but will require more programming

    \item An SDK such as Claude's Agent SDK\cite{claude-agent-sdk}, which strikes a balance between the two above at the cost of vendor lock-in.
\end{enumerate}

Each implementation must:

\begin{itemize}
  \item Use the same frozen specification,
  \item Be developed independently (fresh session),
  \item Be runnable against the test suite.
\end{itemize}

\textbf{Hint}: I picked a large enough $N$ that you should automate this with an agentic approach, or at minimum a wrapper around a series of LLM calls.

 \begin{tcolorbox}[colback=yellow!10,colframe=black,boxrule=0.5pt]
  \textbf{Deliverable:} Give your derivation of the $k$oo$N$ system failure probability.
  \end{tcolorbox} 

\subsubsection{Step 2: Measure Failure Behavior}

We will release the full test suite on Brightspace.

For each input in the test suite:

\begin{itemize}
  \item Record which versions produce incorrect output.
  \item Estimate:
  \begin{itemize}
    \item Marginal failure probability $p$,
    \item Pairwise joint failure probabilities,
    \item Empirical distribution of the number of failures per input.
  \end{itemize}
\end{itemize}

 \begin{tcolorbox}[colback=yellow!10,colframe=black,boxrule=0.5pt]
  \textbf{Deliverable:} Provide your analysis, illustrated with at least one table or figure.
  \end{tcolorbox} 

\subsubsection{Step 3: Compare to the $k$-out-of-$N$ Model}

Using your derived formula from the Theory section:

\begin{itemize}
  \item Compute the predicted system failure probability under the assumption of independent failures.
  \item Compute the observed system failure rate under $k$-out-of-$N$ voting.
  \item Compare the predicted and observed values.
\end{itemize}

At this stage, you are performing a \emph{descriptive model comparison}, not a formal statistical hypothesis test.
Your goal is to assess whether the empirical behavior appears consistent with the binomial independence model, and to characterize any discrepancies you observe.

Discuss:

\begin{itemize}
  \item Whether the observed joint behavior is qualitatively consistent with independence.
  \item Whether voting improves reliability relative to a single version.
  \item How strongly the observed behavior deviates from the binomial prediction.
  \item Plausible explanations for any deviation.
\end{itemize}

You are not required to compute $p$-values or conduct formal hypothesis tests here.
If you wish to perform a rigorous statistical test of independence, see the optional subsection below.

\subsubsection{Optional: Statistical Analysis in the Style of Knight \& Leveson}

In the required portion of this lab, you compared observed reliability to the prediction of the $k$-out-of-$N$ model under independence.
That comparison is descriptive: you assessed whether the empirical data appear to conform to the binomial model, but you did not perform a formal statistical hypothesis test.

Knight and Leveson (1990), by contrast, conducted formal statistical analyses to evaluate whether observed joint failure behavior was consistent with independence.
Replicate their statistical framing, including:

\begin{itemize}
  \item Testing whether observed joint failure rates significantly exceed $p^2$ (or the appropriate independence prediction for the $k$oo$N$ scenario),
  \item Performing formal hypothesis tests of independence,
  \item Estimating correlation coefficients between version failure indicators,
  \item Computing confidence intervals for marginal and joint failure probabilities.
\end{itemize}

If you choose this option, clearly state:

\begin{itemize}
  \item Your null hypothesis ($H_0$),
  \item The statistical test used,
  \item Any assumptions required,
  \item Your interpretation of the result.
\end{itemize}

\begin{tcolorbox}[colback=yellow!10,colframe=black,boxrule=0.5pt]
\textbf{Deliverable:} Provide your statistical analysis and interpretation.
\end{tcolorbox}

This deliverable is entirely optional, but instructive.
You are welcome to use ChatGPT or other tools to assist with statistical calculations, provided that you understand and clearly explain the methodology used.
You may complete any part of it that you see fit.

\subsection{Deliverables}

Provide your empirical results and a written interpretation (1–2 pages) addressing whether the data conform to the $k$oo$N$ independence model.
You may place your solutions in the colorboxes above, or in a new section here as you like.

\subsection{A Note on Interpreting ``No Failures Observed''}

Design redundancy is most suitable when single-version reliability cannot be assured, or when residual risk remains unacceptably high even after best-effort assurance activities.
If single-version reliability is already extremely high on the input distribution represented by our test suite, then an $N$-version experiment may observe few or no failures, and therefore may have limited ability to evaluate failure independence empirically.

To see why, suppose a single generated implementation fails on approximately $1$ in $10^6$ inputs (\ie $p = 10^{-6}$).
Our controlled test suite contains $T = 10{,}000$ inputs.
The expected number of failures for a single implementation on this suite is:

\[
\mathbb{E}[\text{single-version failures}] = Tp = 10{,}000 \cdot 10^{-6} = 0.01.
\]

Thus, even if a given implementation is imperfect, we should expect to observe zero failures most of the time.

Under the independence model, a 1oo2 voter fails only when both versions fail on the same input, which occurs with probability $p^2$.
The expected number of 1oo2 system failures across the test suite is therefore:

\[
\mathbb{E}[\text{1oo2 system failures}] = T p^2 = 10{,}000 \cdot (10^{-6})^2 = 10^{-8}.
\]

This value is effectively zero.
In such a regime, the absence of observed failures is not strong evidence that failures are independent, nor that redundancy is unnecessary; it primarily indicates that the experiment, as instantiated by the available test suite, has limited statistical power to observe rare events.

If you observe very few (or zero) failures, you should report this outcome explicitly and discuss what it implies---and does not imply---about (i) single-version reliability, (ii) failure independence, and (iii) the value of design diversity for this task.
